% Options for packages loaded elsewhere
\PassOptionsToPackage{unicode}{hyperref}
\PassOptionsToPackage{hyphens}{url}
\documentclass[
  ignorenonframetext,
]{beamer}
\newif\ifbibliography
\usepackage{pgfpages}
\setbeamertemplate{caption}[numbered]
\setbeamertemplate{caption label separator}{: }
\setbeamercolor{caption name}{fg=normal text.fg}
\beamertemplatenavigationsymbolsempty
% remove section numbering
\setbeamertemplate{part page}{
  \centering
  \begin{beamercolorbox}[sep=16pt,center]{part title}
    \usebeamerfont{part title}\insertpart\par
  \end{beamercolorbox}
}
\setbeamertemplate{section page}{
  \centering
  \begin{beamercolorbox}[sep=12pt,center]{section title}
    \usebeamerfont{section title}\insertsection\par
  \end{beamercolorbox}
}
\setbeamertemplate{subsection page}{
  \centering
  \begin{beamercolorbox}[sep=8pt,center]{subsection title}
    \usebeamerfont{subsection title}\insertsubsection\par
  \end{beamercolorbox}
}
% Prevent slide breaks in the middle of a paragraph
\widowpenalties 1 10000
\raggedbottom
\AtBeginPart{
  \frame{\partpage}
}
\AtBeginSection{
  \ifbibliography
  \else
    \frame{\sectionpage}
  \fi
}
\AtBeginSubsection{
  \frame{\subsectionpage}
}
\usepackage{iftex}
\ifPDFTeX
  \usepackage[T1]{fontenc}
  \usepackage[utf8]{inputenc}
  \usepackage{textcomp} % provide euro and other symbols
\else % if luatex or xetex
  \usepackage{unicode-math} % this also loads fontspec
  \defaultfontfeatures{Scale=MatchLowercase}
  \defaultfontfeatures[\rmfamily]{Ligatures=TeX,Scale=1}
\fi
\usepackage{lmodern}
\ifPDFTeX\else
  % xetex/luatex font selection
\fi
% Use upquote if available, for straight quotes in verbatim environments
\IfFileExists{upquote.sty}{\usepackage{upquote}}{}
\IfFileExists{microtype.sty}{% use microtype if available
  \usepackage[]{microtype}
  \UseMicrotypeSet[protrusion]{basicmath} % disable protrusion for tt fonts
}{}
\makeatletter
\@ifundefined{KOMAClassName}{% if non-KOMA class
  \IfFileExists{parskip.sty}{%
    \usepackage{parskip}
  }{% else
    \setlength{\parindent}{0pt}
    \setlength{\parskip}{6pt plus 2pt minus 1pt}}
}{% if KOMA class
  \KOMAoptions{parskip=half}}
\makeatother
\usepackage{longtable,booktabs,array}
\newcounter{none} % for unnumbered tables
\usepackage{calc} % for calculating minipage widths
\usepackage{caption}
% Make caption package work with longtable
\makeatletter
\def\fnum@table{\tablename~\thetable}
\makeatother
\usepackage{graphicx}
\makeatletter
\newsavebox\pandoc@box
\newcommand*\pandocbounded[1]{% scales image to fit in text height/width
  \sbox\pandoc@box{#1}%
  \Gscale@div\@tempa{\textheight}{\dimexpr\ht\pandoc@box+\dp\pandoc@box\relax}%
  \Gscale@div\@tempb{\linewidth}{\wd\pandoc@box}%
  \ifdim\@tempb\p@<\@tempa\p@\let\@tempa\@tempb\fi% select the smaller of both
  \ifdim\@tempa\p@<\p@\scalebox{\@tempa}{\usebox\pandoc@box}%
  \else\usebox{\pandoc@box}%
  \fi%
}
% Set default figure placement to htbp
\def\fps@figure{htbp}
\makeatother
\setlength{\emergencystretch}{3em} % prevent overfull lines
\providecommand{\tightlist}{%
  \setlength{\itemsep}{0pt}\setlength{\parskip}{0pt}}
\usepackage{bookmark}
\IfFileExists{xurl.sty}{\usepackage{xurl}}{} % add URL line breaks if available
\urlstyle{same}
\hypersetup{
  hidelinks,
  pdfcreator={LaTeX via pandoc}}

\author{\texorpdfstring{}{}}
\date{}

\begin{document}

\section{Case Systems: From Pāṇinian Kāraka to Cross-Linguistic
Alignment Typology}\label{sec:case-systems}

\begin{frame}{Five Analytical Traditions That Shaped the Modern Theory
of Case}
\protect\phantomsection\label{five-analytical-traditions-that-shaped-the-modern-theory-of-case}
\begin{block}{The Pāṇinian Kāraka Framework}
\protect\phantomsection\label{the-pux101ux1e47inian-kux101raka-framework}
The formal study of grammatical case originates with Pāṇini (circa 4th
century BCE), whose \emph{Aṣṭādhyāyī}---a grammar of approximately 4,000
rules that predates Euclid---formalized the \emph{kāraka} theory. This
framework first classified semantic roles---agent, patient, instrument,
and others---as deep relational functions linking verbs to their
arguments, abstracting away from surface morphosyntactic inflections.
Etymologically meaning ``that which brings about'' an action, the kāraka
system establishes a rigorous mapping from conceptual predication to
phonological realization {[}-@jha2021sanskrit; -@kak1987paninian{]}.
\end{block}

\begin{block}{Jakobson's Structural Features and the Prague School}
\protect\phantomsection\label{jakobsons-structural-features-and-the-prague-school}
This profound semantic foundation lay predominantly dormant until
structural linguists resurrected it in the mid-20th century. Roman
Jakobson's \emph{Morphologic Inquiry into Slavic Declension} (1958)
decomposed grammatical cases into binary distinctive features (e.g.,
{[}±directional{]}, {[}±peripheral{]}) {[}-@jakobson1958morphologic{]}.
By treating case oppositions analogously to phonological features,
Jakobson shifted the field from Pāṇini's holistic roles to a
componential analysis that exposes deep relational hierarchies and
markedness. Concurrently, Louis Hjelmslev's \emph{La catégorie des cas}
{[}-@hjelmslev1935categorie{]} reinforced this functionalist tradition
by formalizing case as a purely relational category within a dynamic
semiotic network.
\end{block}

\begin{block}{Fillmore's Deep Case and Generative Roots}
\protect\phantomsection\label{fillmores-deep-case-and-generative-roots}
Charles Fillmore's seminal ``The Case for Case'' (1968) built directly
upon this structuralist lineage, explicitly translating these concepts
into the burgeoning generative linguistics framework. Fillmore proposed
\emph{deep cases} (e.g., Agentive, Objective, Dative) as universal
semantic primitives underlying surface syntax. Crucially, he argued that
verbs assign underlying \emph{case frames} which subsequently generate
surface structures {[}-@fillmore1968case{]}. This evolution transformed
Pāṇini's kāraka into deep relational networks, permanently prioritizing
universal semantics over language-specific morphology.
\end{block}

\begin{block}{Mel'čuk's Meaning-Text Theory (MTT)}
\protect\phantomsection\label{melux10duks-meaning-text-theory-mtt}
A distinct but parallel formalization emerged with Aleksandr Žolkovskij
and Igor Mel'čuk's Meaning-Text Theory (MTT) {[}@zolkovskij1965possible;
@melcuk1981meaning; @melcuk1988dependency{]}. Developed initially in
Moscow, MTT models natural language as a rigorous, many-to-many
correspondence between meanings (deep semantic representations) and
texts (surface-phonological representations). The transition from
meaning to text unfolds across a multi-level synthesis process: Deep
Semantics, Deep-Syntactic, Surface-Syntactic, Morphological, and
Phonological. At the deep-syntactic level, meaning is represented via
\emph{dependency trees} linking lexical units through dependency
relations. The nodes are populated by \emph{actants}---semantic roles
closely aligning with thematic grids but strictly language-specific in
their surface realization.

Crucially, MTT establishes that grammatical case is not directly
semantic, but rather a final surface-morphological phenomenon governed
by syntactic linearization and dependency constraints. \emph{Lexical
functions} map actants to surface structures, explicitly demonstrating
that morphosyntactic inflection serves merely as the end-stage formal
realization of deep semantic relations. By monotonically mapping meaning
to text via hierarchical dependency graphs, MTT provided an exhaustive
synthesis that presages our modern categorical formalizations mapping
conceptual structures to syntactic types.
\end{block}

\begin{block}{Dowty's Proto-Roles and Graded Topologies}
\protect\phantomsection\label{dowtys-proto-roles-and-graded-topologies}
Fillmore's deep cases serve as the direct precursors to modern
\emph{thematic role} theory. Dowty {[}-@dowty1991thematic{]} refined
this paradigm by decomposing thematic roles into clusters of sentential
entailments, yielding two \emph{proto-roles}: the Proto-Agent
(characterized by volitional involvement and causation) and the
Proto-Patient (characterized by incremental themes and causal
affectedness). This decomposition proves critical for our categorical
formalization: it replaces discrete, named roles with a \emph{graded}
structural topology where role assignment is a matter of degree rather
than kind. Consequently, morphisms in our cognitive case diagrams carry
continuous weights \(w \in [0,1]\) representing the degree to which a
noun phrase satisfies proto-role entailments---a design choice that
yields the enriched category structure developed formally in
\autoref{sec:enriched-categories}.
\end{block}
\end{frame}

\begin{frame}[fragile]{Four Alignment Systems}
\protect\phantomsection\label{four-alignment-systems}
Contemporary typological work reveals that the world's languages realize
case systems according to a small number of \emph{alignment
types}---systematic patterns governing how the core arguments of
transitive and intransitive clauses are grouped {[}@polinsky2015case;
@blake2001grammatical; @haspelmath2009universality{]}.

\begin{block}{S, A, P}
\protect\phantomsection\label{s-a-p}
The cross-linguistic comparison rests on three primitives:

{\def\LTcaptype{none} % do not increment counter
\begin{longtable}[]{@{}
  >{\centering\arraybackslash}p{(\linewidth - 4\tabcolsep) * \real{0.3571}}
  >{\raggedright\arraybackslash}p{(\linewidth - 4\tabcolsep) * \real{0.3571}}
  >{\raggedright\arraybackslash}p{(\linewidth - 4\tabcolsep) * \real{0.2857}}@{}}
\toprule\noalign{}
\begin{minipage}[b]{\linewidth}\centering
Symbol
\end{minipage} & \begin{minipage}[b]{\linewidth}\raggedright
Role
\end{minipage} & \begin{minipage}[b]{\linewidth}\raggedright
Definition
\end{minipage} \\
\midrule\noalign{}
\endhead
\textbf{S} & Sole argument of intransitive & ``The child
\textbf{sleeps}'' \\
\textbf{A} & Agent-like argument of transitive & ``\textbf{The child}
broke the vase'' \\
\textbf{P} & Patient-like argument of transitive & ``The child broke
\textbf{the vase}'' \\
\bottomrule\noalign{}
\end{longtable}
}
\end{block}

\begin{block}{Accusative, Ergative, Active-Stative, Fluid-S: Four Ways
to Group Intransitive Subjects}
\protect\phantomsection\label{accusative-ergative-active-stative-fluid-s-four-ways-to-group-intransitive-subjects}
The key insight from typological research is that languages differ in
how they \emph{group} these three roles for purposes of case marking,
agreement, and other grammatical processes:

{\def\LTcaptype{none} % do not increment counter
\begin{longtable}[]{@{}
  >{\raggedright\arraybackslash}p{(\linewidth - 4\tabcolsep) * \real{0.3333}}
  >{\raggedright\arraybackslash}p{(\linewidth - 4\tabcolsep) * \real{0.3333}}
  >{\raggedright\arraybackslash}p{(\linewidth - 4\tabcolsep) * \real{0.3333}}@{}}
\toprule\noalign{}
\begin{minipage}[b]{\linewidth}\raggedright
Alignment
\end{minipage} & \begin{minipage}[b]{\linewidth}\raggedright
Grouping
\end{minipage} & \begin{minipage}[b]{\linewidth}\raggedright
Exemplar Languages
\end{minipage} \\
\midrule\noalign{}
\endhead
\textbf{Nominative--Accusative} & S = A \(\neq\) P & English, Latin,
Finnish, Russian \\
\textbf{Ergative--Absolutive} & S = P \(\neq\) A & Basque, Dyirbal,
Georgian (partly) \\
\textbf{Active--Stative} & S splits by agentivity & Lakhota, Guaraní,
Eastern Pomo \\
\textbf{Tripartite} & S \(\neq\) A \(\neq\) P & Nez Perce, some
Australian languages \\
\textbf{Fluid-S} & S marking varies by context & Bats (NE Caucasian),
Acehnese \\
\bottomrule\noalign{}
\end{longtable}
}

\textbf{Fluid-S and Context-Dependent Functors.} In Bats
(Nakh-Daghestanian), the intransitive subject of a single verb surfaces
in different cases strictly depending on the speaker's internal
construal of agentive volition. For example, the verb \emph{fall}
assigns an absolutive S when the action is accidental (\emph{The
child-ABS fell}), but assigns an ergative S when the action is
volitional (\emph{The child-ERG fell {[}on purpose{]}}).

Categorically, we model Fluid-S as a \textbf{context-dependent functor}
\(F_\theta: \mathcal{U} \to \mathcal{L}\) parameterized by a continuous
volition feature \(\theta \in [0,1]\). \autoref{fig:fluid-s} visualizes
the resulting volition landscape: case categorization boundaries shift
dynamically as a direct function of the agent's internal construal,
satisfying naturality only up to a probabilistic reparameterization of
\(\theta\).

\begin{figure}
\centering
\pandocbounded{\includegraphics[keepaspectratio,alt={Fluid-S alignment is a continuous mapping, not a binary switch. The context-dependent functor F\_\textbackslash theta: \textbackslash mathcal\{U\} \textbackslash to \textbackslash mathcal\{L\} is rendered as a 2D decision surface with volitional control \textbackslash theta \textbackslash in {[}0,1{]} on the x-axis and proto-agentivity {[}@dowty1991thematic{]} on the y-axis. Color intensity represents P(\textbackslash text\{ERG\} \textbackslash mid \textbackslash theta, \textbackslash text\{agentivity\}), computed via a logistic boundary. Nakh-Daghestanian verb exemplars (Bats language) are overlaid at their typologically attested coordinates: low-volition actions (sneeze, accidental fall) cluster in the ABS region; high-volition actions (jump, fight) occupy the ERG region. The dashed curve marks the functor decision boundary where F\_\textbackslash theta(S) transitions from ABS to ERG. Generated programmatically from the FluidSFunctor class.}]{output/figures/fluid_s_volition_landscape.png}}
\caption{Fluid-S alignment is a continuous mapping, not a binary switch.
The context-dependent functor \(F_\theta: \mathcal{U} \to \mathcal{L}\)
is rendered as a 2D decision surface with volitional control
\(\theta \in [0,1]\) on the x-axis and proto-agentivity
{[}@dowty1991thematic{]} on the y-axis. Color intensity represents
\(P(\text{ERG} \mid \theta, \text{agentivity})\), computed via a
logistic boundary. Nakh-Daghestanian verb exemplars (Bats language) are
overlaid at their typologically attested coordinates: low-volition
actions (sneeze, accidental fall) cluster in the ABS region;
high-volition actions (jump, fight) occupy the ERG region. The dashed
curve marks the functor decision boundary where \(F_\theta(S)\)
transitions from ABS to ERG. Generated programmatically from the
\protect\texttt{FluidSFunctor} class.}\label{fig:fluid-s}
\end{figure}

\textbf{Synthetic Case-Role Algebra.} We introduce Synthetic Case-Role
Algebra as a novel, computational upgrade to the Dowty-style proto-role
framework. Where Dowty modeled proto-roles as static clusters of
entailments {[}-@dowty1991thematic{]}, we formalize them as
\textbf{objects in a \([0,1]\)-enriched monoidal category} with tensor
product over role compositions. This advancement enables purely
algebraic manipulation of semantic roles: composition, weighting, and
transformation of roles proceed through functorial operations
representing complex event structures such as causativization, serial
verb constructions, and argument-structure alternations. Crucially, this
enriched structure demonstrates that ``case assignment'' is not a
discrete binary choice but a vector-valued expectation in continuous
case space---providing the mathematical bridge between the symbolic
traditions of formal grammar and the statistical representations of
modern neural language models (\autoref{sec:categorical-semantics}).

Claassen {[}-@claassen2019alignment{]} provides a comprehensive survey
of the explanatory frameworks proposed for alignment diversity, arguing
that no single factor (processing efficiency, disambiguation, discourse
pragmatics) suffices---a conclusion that motivates our multi-dimensional
categorical formalization. Wu {[}-@wu2024amis{]} offers a detailed case
study of Amis (Austronesian), demonstrating how verb classification,
case marking, and grammatical relations interact in a language that
defies simple alignment classification.

Beyond the three core arguments, languages distinguish a rich inventory
of oblique cases. Our formalization follows the CEREBRUM framework
{[}@friedman2024cerebrum{]} in adopting eight fundamental cases:

{\def\LTcaptype{none} % do not increment counter
\begin{longtable}[]{@{}
  >{\raggedright\arraybackslash}p{(\linewidth - 6\tabcolsep) * \real{0.2353}}
  >{\centering\arraybackslash}p{(\linewidth - 6\tabcolsep) * \real{0.2941}}
  >{\raggedright\arraybackslash}p{(\linewidth - 6\tabcolsep) * \real{0.2353}}
  >{\raggedright\arraybackslash}p{(\linewidth - 6\tabcolsep) * \real{0.2353}}@{}}
\toprule\noalign{}
\begin{minipage}[b]{\linewidth}\raggedright
Case
\end{minipage} & \begin{minipage}[b]{\linewidth}\centering
Abbreviation
\end{minipage} & \begin{minipage}[b]{\linewidth}\raggedright
Semantic Core
\end{minipage} & \begin{minipage}[b]{\linewidth}\raggedright
Syntactic Prototype
\end{minipage} \\
\midrule\noalign{}
\endhead
Nominative & NOM & Agent / experiencer & Intransitive subject,
transitive agent \\
Accusative & ACC & Patient / theme & Direct object, incremental theme \\
Genitive & GEN & Possessor / source & Possessive modifier, partitive \\
Dative & DAT & Recipient / goal & Indirect object, beneficiary \\
Instrumental & INS & Instrument / means & Adverbial of means \\
Locative & LOC & Location / context & Spatial/temporal ground \\
Ablative & ABL & Origin / cause & Source of motion, causal adjunct \\
Vocative & VOC & Addressee & Direct address \\
\bottomrule\noalign{}
\end{longtable}
}

While historically used merely to diagram sentences, this exact
eight-role inventory is what enables the \textbf{Categorical AI
Protocol} introduced in \autoref{sec:ai-implications}. By rigidly
mapping artificial agent capabilities to corresponding grammatical
cases---e.g., treating an API as strictly \texttt{INS}, passive data
strictly as \texttt{ACC}, and system context rigidly as
\texttt{LOC}---the cognitive case diagram enforces computational
boundary constraints that natively repel prompt injection attacks
(\autoref{sec:cognitive-security}).
\end{block}
\end{frame}

\end{document}
